plexe du théorème de Riemann--Roch présenté ici, ou de variantes convenables du théorème de l'index d'Atiyah--Singer. Quelques indications dans ce sens seront données dans Exp.\ II. Malheureusement, il manque encore à l'heure actuelle un énoncé, même heuristique, qui engloberait ces deux théorèmes (dont le premier pour l'instant reste conjectural). Il manque apparemment une idée nouvelle, comme il en manque aussi en Géométrie Algébrique pour parvenir à une démonstration du théorème de Riemann--Roch en dehors d'hypothèses projectives (cf.\ Exp.\ XIV, n\textsuperscript{o}~2).

Nous n'aurons à faire nul usage dans ce séminaire de la \underline{théorie locale} des intersections, exposée dans J.P. Serre, Algèbre Locale. Multiplicités (Lecture Notes in Mathematics n\textsuperscript{o}~11, Springer), et nous contenterons simplement de signaler ici que ce cours de Serre a eu une influence évidente sur la genèse de la théorie en 1957. Nous n'utiliserons pas non plus la théorie de l'\underline{anneau de Chow} développée dans Séminaire Chevalley 1958, exposés 2 et 3 (École Normale Supérieure). Cette théorie est liée de façon essentielle à des hypothèses de non singularité, alors que le but de notre séminaire est au contraire de développer une théorie des intersections sur les schémas généraux (et même les topos localement annelés généraux) voir à ce sujet les commentaires dans Exp.\ XIV n\textsuperscript{o}~4 et n\textsuperscript{o}~8, donnant les relations entre notre théorie et celle de l'anneau de Chow, et posant la question d'une généralisation de cette dernière. Nous pouvons dire que le Séminaire présente une théorie des intersections cohérente et se suffisant à elle même, mais nullement exhaustive des différents points de vue utiles (voire indispensables) en théorie des intersections, et qu'il convient par suite de le compléter par les exposés cités de Serre et de Chevalley.

Nous avons joint au séminaire (en Appendice à Exp.\ 0) le rapport Grothendieck de 1957 «~classes de Faisceaux et théorème de Riemann--Roch~», qui avait servi de base au séminaire de Borel et Serre à Princeton la même année, ainsi qu'à leur article déjà cité. Ce rapport esquisse deux démonstrations du théorème de Riemann--Roch pour les variétés quasi-projectives non singulières, dont la première, valable pour le moment seulement en caractéristique nulle, mais donnant en revanche un résultat un peu plus précis dans le cas d'une immersion, ne figure pas encore dans la littérature (sans exclure le travail de Borel--Serre ni le présent séminaire). La lecture de ce rapport ne présuppose bien entendu celle d'aucun autre exposé du séminaire, et peut même servir d'introduction à l'étude de ce dernier au même titre que l'exposé 0, surtout pour un lecteur qui ne serait pas encore familier avec Borel--Serre. La démonstration à laquelle nous venons de faire allusion s'applique d'ailleurs essentiellement telle quelle au cas 
